\ProvidesFile{ArgDispatcher.tex}[Автоматическая диспетчеризация аргументов]


\subsection{Автоматическая диспетчеризация аргументов}
\label{section::ArgDispatcher}

\subsubsection{Что мы хотим}

Предположим, что у нас есть заранее известный список переменных разного типа:
\[
\verb"T1 t1",\; \verb"T2 t2",\; \ldots,\; \verb"Tn tn"
\]
И у нас есть функция \verb"g", у которой аргументы имеют разный тип и эти типы лежат среди указанных выше.
Например,
\begin{cppcode}[\nonumbers]
int g(T4 x, const T1& y, T2& z)
\end{cppcode}
И теперь мы хотим подставить переменные \verb"t1",\ldots,\verb"tn" в функцию \verb"g" автоматически.
То есть хотим написать вызов вида:
\begin{cppcode}[\nonumbers]
int res = auto_call(g, t1, ..., tn);
\end{cppcode}
чтобы он был эквивалентен вызову
\begin{cppcode}[\nonumbers]
int res = g(t4, t1, t2);
\end{cppcode}
То есть мы хотим идентифицировать аргументы по их типам.
Я думаю, что можно придумать другие более хитрые способы диспетчеризации, но для начала давайте разберемся с этим.
При этом мы хотим сообщать об ошибках в следующих случаях:
\begin{enumerate}
\item Если список переменных \verb"t1",\ldots,\verb"tn" содержит переменные одинакового типа (возможно отличающиеся \verb"const/volatile").
В этом случае мы не сможем определить какую переменную мы хотим подставлять только по ее типу.

\item Если список аргументов \verb"g" содержит аргументы одинакового типа.
Здесь мы не отличаем тип от ссылки на тип, так же не отличаем наличие модификаторов \verb"const/volatile".
Например, функция
\begin{cppcode}[\nonumbers]
int g(int x, const int& y, int& z)
\end{cppcode}
для нас имеет три одинаковых типа в аргументах -- \verb"int".
Технически мы можем считать, что все эти ситуации предполагают, что мы подставляем одну и ту же переменную во все три аргумента.
И это даже не очень плохо.

Так же можно было бы заставить пользователя явно каставать аргументы с помощью \verb"std::move", например так:
\begin{cppcode}[\nonumbers]
int x = 1;
int g(1, static_cast<const int&>(x), x);
\end{cppcode}
Но вот тут поди догадайся \verb"x" биндится на аргумент типа \verb"int", \verb"const int&" или \verb"int&".
Так что этот подход не годится в принципе.
Потому я предпочитаю жесткое ограничение -- никаких повторений типов даже на уровне ссылок.

\item Если тип аргумента \verb"g" не содержится среди типов \verb"T1",\ldots,\verb"Tn".
В этом случае мы просто не можем выбрать среди \verb"t1",\ldots,\verb"tn" аргументы для \verb"g".
Можно было бы рассматривать возможность неявной конвертации, но это может привести к неожиданным проблемам, когда компилятор не сможет выбрать нужный тип.
\end{enumerate}

Давайте сразу скажу про одну проблему, которая у нас возникает.
В языке возможно перегружать функции.
Например
\begin{cppcode}[\nonumbers]
int g(int, char, double);
int g(int, float);
\end{cppcode}
То мы просто так не сможем вызвать
\begin{cppcode}[\nonumbers]
int res = auto_call(g, t1, ... tn);
\end{cppcode}
Потому что компилятор не знает, какую из функций мы имеем в виду.
Но эта проблема решается с помощью выбора сигнатуры явно с помощью \verb"sig" шаблона разобранного в разделе~\ref{section::FunctionList}.%
\footnote{Сама имплементация начинается в разделе ~\ref{section::Sig}.}
\begin{cppcode}[\nonumbers]
int res = auto_call(sig<int(int, float)>(g), t1, ... tn);        // call the second function
int res = auto_call(sig<int(int, char, double)>(g), t1, ... tn); // call the first function
\end{cppcode}

\subsubsection{Определение сигнатуры}

Прежде чем решать саму задачу, нам надо научиться решать вспомогательную задачу, которая является в некотором смысле обратной к той, которую мы решали в разделе~\ref{section::FunctionList}.
Нам нужно для каждой функции определять ее сигнатуру.
Кроме того, мы не обязательно можем вызывать функцию.
первый аргумент в \verb"auto_call" может быть объектом с перегруженным оператором \verb"operator()".
Потому мы должны определять сигнатуру либо этого оператора, либо сигнатуру функции.

\paragraph{Сигнатура функции}

Для определения сигнатуры функции сделаем следующий вспомогательный шаблон
\begin{cppcode}
namespace detail {
template<class F>
struct SigOfImpl {
  static_assert(False<F>,
                "The argument of SigOf MUST be a function with no overloads or "
                "a callable object with a unique operator()!");
};

template<class R, class... Args>
struct SigOfImpl<R(Args...)> {
  using type = TL::List<R, Args...>;
};

template<class R, class... Args>
struct SigOfImpl<R (*)(Args...)> {
  using type = TL::List<R, Args...>;
};
} // namespace detail
\end{cppcode}
Обратите внимание, что мы не рассматриваем методы классов, потому что для них нужен будет еще один вспомогательный аргумент типа класса.
В целом можно было бы рассматривать и вызов метода класса, но я не буду этого делать тут.

\paragraph{Добавляем \texttt{callable objects}}

Теперь добавим объекты, у которых есть оператор \verb"operator()".
У нас будет одно жесткое ограничение -- этот оператор НЕ должен быть перегружен.
Либо надо будет сделать аналог \verb"sig" для выбора конкретной перегрузки оператора \verb"operator()".
Я этого делать не буду.
Я добавлю вспомогательный аргумент, как флаг, что мы рассматриваем специальный случай объекта.
\begin{cppcode}
namespace detail {
template<class F, class = void>
struct SigOfImpl {
  static_assert(False<F>,
                "The argument of SigOf MUST be a function with no overloads or "
                "a callable object with a unique operator()!");
};

template<class R, class... Args>
struct SigOfImpl<R(Args...), void> {
  using type = TL::List<R, Args...>;
};

template<class R, class... Args>
struct SigOfImpl<R (*)(Args...), void> {
  using type = TL::List<R, Args...>;
};

template<class T, class R, class... Args>
struct SigOfImpl<R (T::*)(Args...), void> {
  using type = TL::List<R, Args...>;
};

template<class T, class R, class... Args>
struct SigOfImpl<R (T::*)(Args...) const, void> {
  using type = TL::List<R, Args...>;
};

template<class F>
struct SigOfImpl<F, ForgetV<&F::operator()>> {
  using type = typename SigOfImpl<decltype(&F::operator())>::type;
};
} // namespace detail
\end{cppcode}
Обратите внимание, что во-первых, я добавил \verb"SigOfImpl" для случая методов класса.
Это связано с тем, что я буду определять сигнатуру \verb"T::operator()", который не является глобальной функцией.
А во-вторых, посмотрите, как написана специализация для оператора \verb"operator()".
В последней специализации, мы явно обращаемся к нему по адресу и забываем что мы сделали с помощью шаблона \verb"ForgetV", который подставляет \verb"void" вместо себя.%
\footnote{Как такие трюки работают и что такое \verb"ForgetV" я подробно рассказывал в разделе~\ref{section::LogicChecks}.}
Такой трюк позволяет провалиться в этот случай при наличия единственной версии оператора \verb"operator()" в классе.
А дальше я просто определяю его сигнатуру, игнорируя тип класса.

\paragraph{Пользовательский интерфейс}

Теперь можно было бы сигнатуру вычленить так:
\begin{cppcode}[\nonumbers]
template<class F>
using SigOf = typename detail::SigOfImpl<F>::type;
\end{cppcode}
Однако, этот подход не сработает, если вы передадите ссылку на объект типа \verb"F" или если вы передадите константный объект.
А потому нам надо почистить \verb"F" от ссылок и констант так:
\begin{cppcode}
namespace detail {
...
template<class F>
using Clear = std::remove_cv_t<std::remove_reference_t<F>>;
} // namespace detail

template<class F>
using SigOf = typename detail::SigOfImpl<detail::Clear<F>>::type;
\end{cppcode}
Так же для удобства можно сделать версию от объекта
\begin{cppcode}[\nonumbers]
template<auto f>
using SigOfF = typename detail::SigOfImpl<decltype(f)>::type;
\end{cppcode}
И теперь можно добавить вспомогательные шаблоны для получения типа возвращаемого значения и списка типов аргументов.
\begin{cppcode}[\nonumbers]
template<class F>
using ReturnOf = TL::Front<SigOf<F>>;

template<auto f>
using ReturnOfF = typename TL::Front_t<SigOfF<f>>::type;

template<class F>
using ArgsOf = TL::PopFront<SigOf<F>>;

template<auto f>
using ArgsOfF = TL::PopFront<SigOfF<f>>;
\end{cppcode}
Все, теперь мы можем определить список аргументов для любой функции так
\begin{cppcode}
int g(int, const int&, int&&);
using ArgList = ArgsOfF<g>;
// ArgList equals List<int, const int&, int&&>
\end{cppcode}

\subsubsection{Имплементация}

Начнем имплементацию с общей канвы.
Мы хотим что-то такое:
\begin{cppcode}
template<class F, class... Args>
constexpr auto auto_call(F&& f, Args&&... args) {
  using ArgList = ArgsOf<F>;
  using List = TL::List<Args...>;
  // map args... to ArgList positions according to their types
  // call f(permute(args));
}
\end{cppcode}
Для выполнения перестановки аргументов в нужном порядке, мы в начале найдем список индексов, чтобы понимать, под каким номером надо брать аргумент из списка \verb"args...".
Прежде всего почистим оба списка от ссылок и \verb"const/volatile".
Для этого заведем вспомогательный шаблон
\begin{cppcode}[\nonumbers]
namespace detail {
template<class List>
using RemoveRefCV =
    TL::Map<std::remove_cv_t, TL::Map<std::remove_reference_t, List>>;
} // namespace detail
\end{cppcode}
Теперь почистим списки аргументов:
\begin{cppcode}[\nonumbers]
using ArgList = ArgsOf<F>;
using List = TL::List<Args...>;
using CleanArgList = detail::RemoveRefCV<ArgList>;
using CleanList = detail::RemoveRefCV<List>;
\end{cppcode}
Теперь для каждого типа из \verb"CleanArgList" найдем его позицию в \verb"CleanList".
Для этого определим вспомогательный шаблон:
\begin{cppcode}
namespace detail {
template<class List>
struct Bind {
  template<class T>
  struct IndexOf : TL::IndexOfFirst_t<List, T> {};
};
} // namespace detail
\end{cppcode}
Здесь \verb"Bind" нужен для запоминания \verb"CleanList" как контекста.
Так же обратите внимание, что \verb"IndexOf" является классом, содержащем индекс как \verb"value" поле, а не числом.
Теперь поиск индексов делается так:
\begin{cppcode}
using IndList = TL::Map<detail::Bind<CleanList>::template IndexOf, CleanArgList>;
using EIndList = EL::FromTList<IndList, EL::List, int>;
\end{cppcode}
Здесь во второй строке мы конвертируем индекс к списку вида
\begin{cppcode}
template<class T, T... Es>
struct List{};
// for example EIndList equals List<int, 4, 1, 5>
\end{cppcode}
Теперь нужно вызвать вспомогательную функцию, которая примет \verb"EIndList" как дополнительный аргумент и добавить проверки на ошибки.
Начнем с проверок на ошибки
\begin{cppcode}
static_assert(TL::Size<CleanList> == TL::Size<TL::NoDuplicates<CleanList>>,
              "The arguments of auto_call MUST be of different type!");
static_assert(TL::Size<CleanArgList> == TL::Size<TL::NoDuplicates<CleanArgList>>,
              "The arguments of the first argument of auto_call MUST be of "
              "different type!");
static_assert(!EL::Contains<EIndList, -1>,
              "The types of the arguments of the first argument of auto_call "
              "MUST be among the types of the arguments of auto_call!");
\end{cppcode}
Можно было бы добавить вспомогательную функцию \verb"hasDuplicates" и проверять условие 
\[
\verb"!hasDuplicates<CleanList>"\text{ или }\verb"!hasDuplicates<CleanArgList>"
\]
но я не заморачивался.
Так же обратите внимание, что отсутствие элемента означает, что индекс равен \verb"-1".
Именно это значение я использую как флаг для проверки в третьем условии.
И последняя строчка будет выглядеть так:
\begin{cppcode}[\nonumbers]
return detail::Call_impl<EIndList>::exec(
    std::forward<F>(f), std::forward_as_tuple(std::forward<Args>(args)...));
\end{cppcode}
Здесь структура \verb"Call_impl" нужна для перехвата контекста, которым будет являться список индексов.
А аргументы я передаю в виде \verb"std::tuple" из ссылок на исходные аргументы.
При этом сохраняется вид ссылки lvalue или rvalue.
\verb"exec" является статической функцией внутри \verb"Call_impl" уже для выполнения нужной работы.
Вот как выглядит имплементация вспомогательного класса:
\begin{cppcode}
namespace detail {
template<class TList>
struct Call_impl {
  /*error messages*/
};

template<template<class U, U...> class TList, int... Es>
struct Call_impl<TList<int, Es...>> {
  template<class F, class T>
  static auto exec(F&& f, T&& data) {
    return std::forward<F>(f)(std::get<Es>(std::forward<T>(data))...);
  }
};
} // namespace detail
\end{cppcode}
В данном случае \verb"data" является элементом \verb"std::tuple" наполненная ссылками.
С помощью \verb"std::get<Es>" мы выбираем из этого \verb"std::tuple" элементы с нужными индексами и распаковываем argument pack с помощью многоточия.
Ни и всю эту радость подставляем в функцию \verb"f" (или в оператор \verb"operator()" для объекта \verb"f").

\paragraph{Пример использования}

Теперь я могу сделать так:
\begin{cppcode}[\nonumbers]
void f(int&& x, char y) {
  std::cout << "f(" << x << ", " << y << ");\n";
}

int x = 1;
auto_call(f, std::move(x), 'c', 1.2); // Args: int, char, double
auto_call(f, 1, 'c', 1.2);            // Args: int, char, double
\end{cppcode}
Программа выведет на печать
\begin{cppcode}[\nonumbers]
f(1, 'c');
f(1, 'c');
\end{cppcode}
Если же я сделаю так:
\begin{cppcode}[\nonumbers]
void g(int& x, int y) {
  std::cout << "g(" << x << ", " << y << ");\n";
}
int x = 1;
auto_call(g, x, 'c', 1.2);
\end{cppcode}
То будет ошибка компиляции, потому что у \verb"g" два аргумента одинакового типа.
Если же попробую сделать так
\begin{cppcode}[\nonumbers]
void h(float x, double y) {
  std::cout << "h(" << x << ", " << y << ");\n";
}
int x = 1;
auto_call(h, x, 'c', 1.2);
\end{cppcode}
То будет ошибка компиляции, потому что среди аргументов нет типа \verb"float", а потому вместо \verb"float" нечего подставить.

\subsubsection{Имплементация для метода класса}

Здесь мы хотим несколько локализовать данные, а именно, мы хотим связать переменные \verb"t1",\ldots,\verb"tn" и функцию \verb"auto_call" вместе.
Для этого можно спрятать переменные внутри вспомогательного класса контекста, а функцию \verb"auto_call" сделать его методов для доступа к контексту.
Так можно использовать разные контексты в одном и том же месте кода.
Мы хотим что-то вроде
\begin{cppcode}
void f(const int& x, char y) {
  std::cout << "f(" << x << ", " << y << ");\n";
}

Dispatcher d(1, 'c', 1.2); // check that arguments are of different type
d.call(f);
\end{cppcode}
Тогда можно сделать следующий класс:
\begin{cppcode}
template<class... Ts>
struct Dispatcher {
  using List = TL::List<Ts...>;
  using CleanList = detail::RemoveRefCV<List>;
  static_assert(TL::Size<CleanList> == TL::Size<TL::NoDuplicates<CleanList>>,
                "The data in Dispatcher MUST be of different type!");

  Dispatcher(Ts&&... args) : data(std::make_tuple(std::move(args)...)) {
  }

  template<class F>
  auto call(F&& f) {
    using ArgList = ArgsOf<F>;
    using CleanArgList = detail::RemoveRefCV<ArgList>;
    using IndList =
        TL::Map<detail::Bind<CleanList>::template IndexOf, CleanArgList>;
    using EIndList = EL::FromTList<IndList, EL::List, int>;
    static_assert(TL::Size<CleanArgList> ==
                      TL::Size<TL::NoDuplicates<CleanArgList>>,
                  "The arguments of the argument of call MUST be of "
                  "different type!");
    static_assert(
        !EL::Contains<EIndList, -1>,
        "The types of the arguments of the argument of call "
        "MUST be among the types of the values stored in Dispatcher!");
    return detail::Call_impl<EIndList>::exec(std::forward<F>(f), data);
  }

  std::tuple<Ts...> data;
};
\end{cppcode}
По сути класс хранит три поля:
\begin{enumerate}
\item Конструктор.
Он в частности позволяет не использовать явно имена типов в конструкторе.
Обратите внимание, что тут не универсальные ссылки, а rvalue ссылки.
То есть конструктор поглощает внутрь себя все данные.

\item Поле с данными \verb"data".
Оно специально открытое, чтобы можно было с ними манипулировать напрямую.

\item Метод \verb"call", в который подставляется функция или callable object.
Диспетчеризация выполняется внутри этого метода, где по сути мы сводим задачу к использованию уже имплементированного \verb"Call_impl" класса.
\end{enumerate}
